\section{Related Work}\label{sec:related}
\paragraph{Tables of bounds.}
Systematic tables of bounds on $\dopt(n,k)$ go back to \citet{helgert1973}. \citet{verhoeff1985} reconstructed and updated that table in a technical report, and journal updates followed in \citet{verhoeff1987} and \citet{brouwer1993}. The online tables maintained by \citet{grasslTables} record the current lower and upper bounds together with construction and derivation chains. For the first four parameter pairs in Table~\ref{tab:history}, the recorded upper-bound chains end at the nonexistence of $[25,8,10]$ codes \citep{ytrehus1990}. They propagate that result through parity extension, shortening, and residual-code bounds. Our results replace these chains by sharper bounds.

\paragraph{The Griesmer bound and small dimensions.}
For fixed $k$ and $d$, the Griesmer bound \citep{griesmer1960} gives a lower bound $g_2(k,d)$ on the length of a binary linear $[n,k,d]$ code. \citet{baumert1973} showed that for each fixed dimension the bound is attained for all sufficiently large $d$, so determining the minimum length for fixed $k$ is a finite problem. The minimum lengths for $k\le7$ were determined by \citet{baumert1973} and \citet{vantilborg1981}, and the case $k=8$ was completed by \citet{bouyukliev2000}. As a result, $\dopt(n,k)$ is known for all $k\le8$ \citep{kurz2021}. In dimension nine, \citet{kurz2021} showed that the $[46,9,20]$ code is unique and excluded $[47,10,20]$ and $[85,9,40]$ codes. Our main theorem concerns dimension ten; its reduction rests on the classification for dimensions seven and eight.

\paragraph{Classification of linear codes.}
Exhaustive classification up to equivalence is the main tool for nonexistence results at moderate lengths. \citet{kaski2006} give a general treatment of classification algorithms for codes and designs. \citet{jaffe2000} classified the optimal binary linear codes of length at most 30. \citet{bouyukliev2001} classified optimal binary linear codes of dimension at most seven; their classification supplies the 172 full-support $[40,7,18]$ classes used here. Further classification and nonexistence results for codes with prescribed minimum distance are given by \citet{feulner2014}. \citet{bouyukliev2011} classified all codes of minimum distance 8 and codimension at most 14, and all codes of minimum distance 10 and codimension at most 18, and excluded $[33,18,8]$ and $[33,14,10]$ codes. \citet{bouyukliev2021} developed classification algorithms based on canonical augmentation and lattice-point enumeration. \citet{bouyuklieva2020} classify codes by extending their residuals; their Proposition~2 states that no $[41,14,14]$ code exists, which is the input for our $(42,16)$ result. \citet{bouyuklieva2025} study the hulls of optimal codes, and their counts provide an independent consistency check on our catalogue. In our computation, coordinate equivalence is tested through graph isomorphism with nauty \citep{mckay2014}.

\paragraph{Weight-distribution methods.}
The MacWilliams identities \citep{macwilliams1963} relate the weight distributions of a code and its dual. Linear-programming arguments based on them are a standard source of upper bounds \citep{macwilliams1977}. \citet{schrijver2005} strengthened such bounds for unrestricted binary codes using semidefinite programming. We use the identities differently: a single nonnegative polynomial forces the existence of a dual codeword of weight four, which converts a nonexistence question into an extension problem over a classified subcode. Residual-code arguments combined with the MacWilliams identities also appear in \citet{kurz2021}.

\paragraph{Computer-assisted nonexistence proofs.}
Large finite searches have settled several nonexistence questions in combinatorics; a prominent example is the nonexistence of a projective plane of order 10 \citep{lam1989}. Such proofs are only as reliable as their enumeration. We therefore separate the completeness of each enumeration from the validity of each rejection, cross-check completeness with two independent implementations, and replay every rejection with an independent checker. The deductive steps, as distinct from these computations, are formally checked. The proof search and verification were organized with agent-based research workflows \citep{yang2026aris,mathrsi}.
